\documentclass[11pt]{article}
\usepackage[margin=1in]{geometry}
\usepackage{amsmath,amssymb,amsthm,booktabs,needspace}
\usepackage[hidelinks]{hyperref}
\setlength{\emergencystretch}{2em}
\newtheorem{theorem}{Theorem}
\newtheorem{lemma}[theorem]{Lemma}
\newtheorem{corollary}[theorem]{Corollary}
\newcommand{\mr}{\operatorname{mr}}
\newcommand{\rank}{\operatorname{rank}}
\title{Finite block actions and four periodic twist families\newline satisfying the meridional rank conjecture}
\author{Internal research working proof}
\date{17 November 2026}
\begin{document}
\maketitle

Let $D_j$ be the explicitly listed thirteen-crossing PD for
$j\in\{3145,3147,3148,3149\}$ in Section~\ref{sec:inputs}.
All four start with $(5,1,6,26)$. Replace that first crossing
by the alternating integer tangle of length $1+6N$, leaving
all other crossings unchanged, to obtain $K_j(N)$.
The precise local PD formula is given in Section~\ref{sec:pd}.

\begin{theorem}\label{thm:families}
For every $N\ge0$ and each of these four choices of $j$,
\[
                     \mr(K_j(N))=b(K_j(N))=4.
\]
Moreover,
\[
\begin{array}{c|c}
j&\det K_j(N)\\\hline
3145,3148&543+1776N\\
3147,3149&537+1680N.
\end{array}
\]
Consequently each parametrized family contains infinitely
many distinct knot types. At $N=0$ the four seeds are the
named knots $13a_{3145},13a_{3147},13a_{3148},13a_{3149}$.
\end{theorem}

The four finite seed proofs are retained separately. Here
their 25-point cover actions are induced to 75-point knot
actions, giving a finite quotient criterion which survives
periodic twisting. The proof below is uniform in $N$;
the twelve checked sample diagrams are input controls,
not an interpolation argument. The result does not assert
that the four families are mutually disjoint, that these
knots were previously unresolved worldwide, or that full
classical MRC is solved. The written proof is unformalized
and externally unreviewed.

\section{A finite block-action lower bound}

Let $Q$ act on $r$ disjoint blocks of equal finite size,
transitively on the blocks. Call an involution an exchange
if it interchanges exactly two blocks by mutually inverse
bijections and fixes every point of the other blocks.
For a chosen block $B$, write $Q_B^B$ for the permutation
group induced on $B$ by its setwise stabilizer.

\begin{lemma}\label{lem:gauge}
If $Q$ is generated by $d$ exchanges, then
\[
                         \rank Q_B^B\le d-r+1.
\]
In particular, if $Q_B^B$ is nonabelian then $d\ge r+1$.
\end{lemma}
\begin{proof}
Make a graph with the blocks as vertices and one edge for
each generating exchange. The graph is connected, since
the generators act transitively on the blocks. Choose a
spanning tree, identify the points of one block with a fixed
set, and extend this identification along the tree exchange
bijections. In these coordinates every tree exchange moves
points between its two blocks without changing their labels.
This simultaneous choice is possible because the tree has
no cycles.

Each of the $d-r+1$ remaining edges carries one permutation
$P_e$ of the common label set; traversing it in the other
direction applies $P_e^{-1}$. Every word returning to $B$
therefore acts on its labels by a word in these $P_e$.
Conversely, traverse the tree from $B$ to one endpoint of
$e$, traverse $e$, and return along the tree. This realizes
$P_e$ on $B$. Thus $Q_B^B=\langle P_e\rangle$, giving the
rank bound. A group with at most one generator is cyclic.
\end{proof}

\begin{corollary}\label{cor:finite}
Suppose a knot group maps onto such a $Q$ with a meridian
mapping to an exchange. If $Q_B^B$ is nonabelian, then
$\mr(K)\ge r+1$.
\end{corollary}
\begin{proof}
Every meridian maps to a conjugate of that exchange.
Conjugation in a block-preserving group retains the exchange
property. Any meridional generating family of the knot
group therefore gives an exchange generating family of $Q$.
Apply Lemma~\ref{lem:gauge}.
\end{proof}

This is a finite version of the simple branched-cover rank
argument. It requires neither the order of $Q$ nor a
classification of its subgroups. It is also independent of
the original affine-symplectic common-form condition.

\section{The four 75-point seed actions}

The retained three-sheeted cover has sheet set $\{0,1,2\}$.
Each source meridian arc $i$ exchanges two sheets $u_i,v_i$.
Its lifted edge has an explicit permutation $P_i$ on 25
points; spanning-tree edges have $P_i=1$. The earlier
certificate supplies all eleven non-tree edge images.
On the 75 points $(s,x)$ define
\[
 X_i(u_i,x)=(v_i,P_i(x)),\qquad
 X_i(v_i,x)=(u_i,P_i^{-1}(x)),
\]
and let $X_i$ fix $(s,x)$ when $s\notin\{u_i,v_i\}$.
Each $X_i$ is an exchange of two 25-point blocks.

All thirteen literal involutive Wirtinger relations
$X_oX_aX_oX_b=1$ are checked on all 75 points. Independently,
all original cover relators already pass on 25 points.
These direct evaluations fix the multiplication convention:
products act rightmost first. The set of $X_i$ generates
the finite group denoted by $Q_j$; its full order is not
needed and has not been computed.

Two explicit words $U,V$ stabilize block $B_0$ and have a
nontrivial commutator there:
\[
\begin{array}{c|cc|c}
j&U&V&[U,V](i)\\\hline
3145&X_2X_5X_1&X_2X_{11}X_1&0\mapsto11\\
3147&X_2X_5X_1&X_2X_{11}X_1&5\mapsto9\\
3148&X_2X_5X_1&X_2X_{12}X_1&0\mapsto3\\
3149&X_2X_5X_1&X_2X_{12}X_1&0\mapsto2.
\end{array}
\]
Here the $X_i$ use the one-based retained source arc indices,
and the moved points are the zero-based labels within $B_0$.
The explicit arrays and the source-index transport are in
\path{periodic_twist_family_results.json}.
The quotient on the three blocks contains the tree
transpositions $(0\,1)$ and $(0\,2)$, so it is transitive.
Corollary~\ref{cor:finite} gives lower four for each seed.

Every listed generator image also occurs on an edge outside
the first crossing. This matters below: preserving the
outside edge labels preserves the whole image $Q_j$, not
just a potentially smaller subgroup.

\section{Periodic twisting preserves this lower bound}

For involutions $x,y$, the Wirtinger rule at an integer
crossing sends its ordered incoming labels to
\[
                    T(x,y)=(y,yxy).
\]
Their product $c=xy$ is unchanged. A direct calculation gives
\[
                 T^2(x,y)=(c^{-1}xc,c^{-1}yc).
\]
If $c^m=1$, then $T^{2m}(x,y)=(x,y)$.
This calculation is a group identity, independent of the
permutation degree or a choice of orientations, since all
meridian labels here are involutions.

For the incoming pair at the selected first crossing in
each seed action, $xy$ has order three. Consequently
\[
                         T^{1+6N}(x,y)=T(x,y)
\]
for every $N\ge0$. The new tangle has exactly the original
four boundary labels. Keep every outside label and extend
the interior labels by $T$. Every new Wirtinger relation
is then satisfied, and every interior label is a conjugate
of $x$ or $y$ in $Q_j$. All original generator images remain
on the outside, so the new image is exactly $Q_j$.
Its block-stabilizer action is still nonabelian.
Corollary~\ref{cor:finite} proves
$\mr(K_j(N))\ge4$ uniformly in $N$.

Periodic finite-quotient labellings of twist regions have
precedent. Joseph--Pongtanapaisan's Section 3 constructs
classical sharp examples using overlapping finite cycles
and periodic twist-box rules \cite{jp}. Their explicit
construction and its hypotheses should receive credit;
no novelty claim for twist periodicity is made here.
The block-stabilizer test above uses the retained cover
witnesses, rather than the cycle-generation capacity in
that source. No membership or nonmembership in its known
families has been established for our seeds.

\section{A uniform four-seed upper certificate}

Delete the first crossing but retain its four boundary
half-edges. The outside Wirtinger arcs are the connected
components of edge labels identified across overpasses at
the other twelve crossings. There are fourteen such arcs.
Seed the outside arcs containing these original edge labels:
\[
\begin{array}{c|c}
j&\text{four representative seed edges}\\\hline
3145&1,2,4,18\\
3147&1,2,4,20\\
3148&1,2,4,20\\
3149&1,2,4,20.
\end{array}
\]
The arcs at incoming ports 0 and 1 of the deleted crossing
are already seeded. Once two arcs at one end of an integer
tangle are colored, one Wirtinger move colors the outgoing
underarc at its first crossing; the overarc continues to
the other outgoing strand. Both incoming arcs at the next
crossing are now colored. Induction colors the entire
tangle of any positive length, and both arcs at its other end.

Treat that induction as a two-input tangle move. In each
source outside diagram, eight further retained crossing
moves color all remaining outside arcs. The full traces,
with crossing indices, known over/under arcs and each new
arc, are stored in the certificate. Independent graph
traversal reconstructs the outside arcs and checks every
step. Thus the same four outside seeds color every expanded
diagram, whatever $N\ge0$ is. At $N=0$ some cut outside arcs
may merge through the original crossing; this can only
reduce the number of distinct seeds.

The published equality between bridge number and minimum
diagram Wirtinger number \cite{wirt} gives $b(K_j(N))\le4$.
Together with the lower bound and $\mr\le b$, this proves
the equality part of Theorem~\ref{thm:families}.

\section{Exact family diagrams and distinctness}\label{sec:pd}

Write the removed crossing as $(a,b,c,d)=(5,1,6,26)$ and
put $n=1+6N$. Choose fresh edge labels and write
\[
 L_0=a,\ R_0=b,\quad L_n=d,\ R_n=c.
\]
For $0<i<n$ choose distinct new $L_i,R_i$, also distinct
from all original edge labels. Replace the first tuple by
the ordered list
\[
                  (L_i,R_i,R_{i+1},L_{i+1}),\qquad 0\le i<n.
\]
For definiteness the producer takes $L_i=26+i$ and
$R_i=26+(n-1)+i$ for $1\le i<n$. This local PD has the
original cyclic boundary order, and for $n=1$ is exactly
the original crossing. An odd number of strand exchanges
has the same endpoint pairing as one crossing. Thus the
result remains a knot. Along either strand the over/under
status alternates in the inserted chain and agrees at the
boundary with the original alternating diagram. The whole
new diagram is alternating.

For an alternating diagram the determinant is the number
of spanning trees of either checkerboard graph
\cite[Lemma 3.14]{stoim}. In the retained coloring the local
replacement makes $n$ parallel copies of the first edge $e$.
Every spanning tree either avoids these edges, contributing
$\tau(G\setminus e)$, or uses exactly one of them. There are
$n\tau(G/e)$ possibilities in the latter case. Hence
\[
             \det K_j(N)=\tau(G\setminus e)+(1+6N)\tau(G/e).
\]
The two exact counts are
\[
\begin{array}{c|rr}
j&\tau(G\setminus e)&\tau(G/e)\\\hline
3145,3148&247&296\\
3147,3149&257&280.
\end{array}
\]
One checker rebuilds regions by corner identifications
and enumerates spanning trees directly. A separate Fox
matrix at $-1$, with rational elimination, independently
checks the determinants of the twelve retained diagrams
with $N=0,1,2$. The uniform parallel-edge argument, not
these samples alone, proves the stated affine formulas.
The slopes are positive, so the determinant invariant
separates all members within each family. This completes
Theorem~\ref{thm:families}.

\section{Literal source PDs and assurance}\label{sec:inputs}

All four PDs have the common seven-tuple prefix
\begin{align*}
 &(5,1,6,26),\ (1,15,2,14),\ (7,2,8,3),\ (3,16,4,17),\\
 &(9,5,10,4),\ (13,6,14,7),\ (15,9,16,8).
\end{align*}
Append the appropriate six-tuple column below, in order:
\begin{center}\small
\begin{tabular}{cccc}
\toprule
$3145$&$3147$&$3148$&$3149$\\\midrule
$(19,10,20,11)$&$(21,11,22,10)$&$(21,11,22,10)$&$(23,10,24,11)$\\
$(11,25,12,24)$&$(11,25,12,24)$&$(11,19,12,18)$&$(11,19,12,18)$\\
$(21,13,22,12)$&$(19,12,20,13)$&$(23,12,24,13)$&$(21,13,22,12)$\\
$(17,23,18,22)$&$(17,22,18,23)$&$(17,22,18,23)$&$(17,23,18,22)$\\
$(23,19,24,18)$&$(23,18,24,19)$&$(19,25,20,24)$&$(19,24,20,25)$\\
$(25,20,26,21)$&$(25,21,26,20)$&$(25,21,26,20)$&$(25,20,26,21)$\\
\bottomrule
\end{tabular}
\end{center}
The earlier provenance audit identifies these literally
with two archived KnotInfo releases and the independent
HT census, allowing the standard one-based label shift.
No approximate manifold recognition is involved.

The frozen plan confines the calculation to these four
seeds and their 26 local crossing/axis templates each.
The first crossing and first axis supply the selected
uniform families. Old campaigns on 293 panel names are
not repeated. The independent verifier checks the finite
source actions pointwise, noncommuting stabilizer words,
outside traces, the twelve literal expanded PDs, all their
meridian relators, four-seed propagation and independent
Fox determinants. Changed meridian images, a false period,
an invalid tangle trace and a wrong determinant slope are
rejected. The source parser is shared only for transport
of the retained arc indices. The new pointwise evaluation,
outside graph traversal, corner-region reconstruction,
tree enumeration and determinant elimination are separate
from the producer's implementations.

The literature/priority audit remains incomplete. In
particular, a non-series-parallel checkerboard graph by
itself is insufficient to exclude an arborescent knot:
the published source explicitly warns about alternating
arborescent links with a nonarborescent alternating diagram
\cite[Proposition 6.3 and Remark 6.4]{stoim}.
No worldwide novelty or prior-family exclusion is inferred
from a missing database field or from this graph shortcut.
The separate sixteen-target prefix is settled within the existing
BJW18 graph-cover MRC family; see
\path{ORIGINAL_SIXTEEN_KNOWN_GRAPH_COVER_FAMILY.tex}. Full classical MRC
remains open. The cited failure of meridional-rank
additivity concerns knotted spheres in four dimensions.

\begin{thebibliography}{9}
\bibitem{wirt} R. Blair, A. Kjuchukova, R. Velazquez and P. Villanueva,
Wirtinger systems of generators of knot groups,
\emph{Communications in Analysis and Geometry} 28 (2020), 243--262,
Theorem 1.3.
\url{https://archive.intlpress.com/site/pub/files/_fulltext/journals/cag/2020/0028/0002/CAG-2020-0028-0002-a002.pdf}.
\bibitem{stoim} A. Stoimenow, Graphs, determinants of knots and
hyperbolic volume, \emph{Pacific Journal of Mathematics} 232 (2007),
423--451. Lemmas 3.14 and 3.21; Proposition 6.3 and Remark 6.4.
\url{https://doi.org/10.2140/pjm.2007.232.423}.
\bibitem{jp} J. Joseph and P. Pongtanapaisan,
Meridional rank and bridge number of knotted 2-spheres,
\emph{Canadian Journal of Mathematics} 77 (2025), 282--299,
published online 27 December 2023.
\url{https://doi.org/10.4153/S0008414X23000883}.
The classical construction is read in Section 3 of the
6 February 2023 preprint, with its then-current title
\emph{Bridge number and meridional rank of knotted surfaces}:
\url{https://arxiv.org/html/2111.09233v2}.
The full published proof has not been compared line by line.
\end{thebibliography}
\end{document}
